% PROVENANCE (section-level)
%   status      : ANALYSIS_DECISION (the classification is a working
%                 judgement, not an approved plan)
%   produced-by : the AnalysisTBD and InProgress markers in sections/*.tex
%   commit      : dd520a5
%   sample      : n/a (bookkeeping)
%   selection   : n/a
%   corrections : n/a
\section{Open Items}
\label{app:open}

This appendix consolidates the open items marked throughout the note.
Its purpose is to make visible which single input would close the most of
them, so that effort can be spent in the order that unblocks the most work.

Every item is marked in place with \AnalysisTBD{} or \InProgress{}; nothing
here is new, and nothing here is a claim about the analysis.

\subsection{Distribution of open items}
\label{app:open:counts}

\begin{table}[htbp]
  \centering
  \caption{Open items by section and by what would close them. ``Artifact''
  means a plot, table or fit that must be produced by running the analysis
  code; ``decision'' means a choice for the analyst; ``external'' means
  information held outside the two repositories.}
  \label{tab:open-items}
  \small
  \setlength{\tabcolsep}{5pt}
  \begin{tabular}{>{\raggedright\arraybackslash}p{0.34\textwidth} c c c c}
    \hline
    Section & Artifact & Decision & External & Total \\
    \hline
    1 Introduction              & 0 & 1 & 0 & 1 \\
    2 Data and simulation       & 4 & 1 & 2 & 7 \\
    3 Reconstruction            & 6 & 4 & 2 & 12 \\
    4 Truth classification      & 3 & 1 & 0 & 4 \\
    5 Classifier                & 5 & 3 & 0 & 8 \\
    6 Corrections               & 1 & 4 & 1 & 6 \\
    7 Generic $B\bar{B}$        & 3 & 1 & 0 & 4 \\
    8 Signal extraction         & 4 & 3 & 0 & 7 \\
    9 Systematics               & 2 & 2 & 0 & 4 \\
    10 Validation, box opening  & 2 & 3 & 0 & 5 \\
    11 Results                  & 1 & 0 & 0 & 1 \\
    12 Summary                  & 0 & 1 & 0 & 1 \\
    Appendix: selections        & 0 & 1 & 0 & 1 \\
    Appendix: categories        & 0 & 1 & 2 & 3 \\
    Appendix: classifier        & 6 & 1 & 0 & 7 \\
    Appendix: fit               & 5 & 1 & 0 & 6 \\
    Appendix: open items        & 0 & 3 & 0 & 3 \\
    \hline
    Total                       & 42 & 31 & 7 & 80 \\
    \hline
  \end{tabular}
\end{table}

The totals are generated from the markers in the source and match them
exactly; the split into the three columns is a working judgement.
\AnalysisTBD{the classification into artifact, decision and external in
Table~\ref{tab:open-items} is a working judgement made while writing, not an
approved plan.  Correct it where the intended route to closing an item
differs}

\subsection{Artifacts that would close the most items}
\label{app:open:artifacts}

\begin{enumerate}
  \item \textbf{Run the truth-category validation on one MC16rd file.}
        Closes the \texttt{D\_mcErrors} $> 512$ question of
        Sec.~\ref{sec:truth:closure}, the yield of the three categories
        excluded from the templates, and part of the composition table of
        Sec.~\ref{sec:truth:composition}.
        It is the cheapest item on this list.
  \item \textbf{Regenerate one workspace from the pinned commit.}
        Replaces the superseded tables of Appendix~\ref{app:fit:stored} and
        settles the channel, sample and modifier description of
        Sec.~\ref{sec:fit:modifiers}, which currently documents an obsolete
        realisation of the model.
  \item \textbf{Produce the per-category yield table} in the signal region and
        the sideband, per run period and channel.
        Closes items in Sec.~\ref{sec:truth}, Sec.~\ref{sec:fit} and
        Sec.~\ref{sec:data}.
  \item \textbf{Produce the classifier diagnostics} of
        Appendix~\ref{app:mva:todo}, in particular the correlation of each
        retained input with $\mmiss$ and $\pdl$ in the fit bins, which is the
        only quantitative evidence that the selection does not sculpt the fit
        observables.
  \item \textbf{Obtain a generic-$B\bar{B}$ fit with an interior minimum.}
        This is the precondition for the systematic of
        Sec.~\ref{sec:syst:bbbar} and for constraining the family
        normalisations in Sec.~\ref{sec:fit:modifiers}.
\end{enumerate}

\subsection{Decisions required}
\label{app:open:decisions}

The following are choices rather than measurements, and each is referenced at
the point in the note where it applies.
\begin{itemize}
  \item Which stored classifier model is nominal
        (Sec.~\ref{sec:mva:training}).
  \item Whether the generic-$B\bar{B}$ normalisations are fixed or constrained
        in the signal-region fit, and how the constraint width is obtained
        (Sec.~\ref{sec:fit:modifiers}).
  \item The treatment of form factors and of branching fractions outside the
        generic-$B\bar{B}$ families (Sec.~\ref{sec:corrections:missing}).
  \item The prefit and postfit closure criteria, and the box-opening plan
        (Sec.~\ref{sec:validation:closure},
        Sec.~\ref{sec:validation:unblinding}).
  \item Whether mode replacement is enabled in the nominal analysis
        (Appendix~\ref{app:categories:replacement}).
  \item Whether the reconstruction-level defects recorded in
        Sec.~\ref{sec:recon:b} are corrected before the next production, and
        what is regenerated if they are.
\end{itemize}

\subsection{Items that invalidate existing material if resolved}
\label{app:open:invalidating}

Two open items are not additive: resolving them changes material already
written.
\begin{itemize}
  \item \textbf{The rest-of-event track mask} (Sec.~\ref{sec:recon:b}).
        Correcting it changes which tracks enter the ROE, and therefore
        $M_{\mathrm{bc}}^{\mathrm{ROE}}$, $\Delta E^{\mathrm{ROE}}$, the ROE
        charge and multiplicity, and $\mmiss$.
        Since $\mmiss$ is a fit observable and the ROE multiplicity enters the
        generic-$B\bar{B}$ tuning objective, every ROE-derived number in this
        note predates the fix.
  \item \textbf{The classifier} (Sec.~\ref{sec:mva}).
        The generic-$B\bar{B}$ tuning region is defined by classifier output,
        so retraining requires re-deriving the family weights.
\end{itemize}
Work that depends on either should be sequenced after them rather than
before.
